% MT033 / answer-en — canonical editable content.
A generative approach models the distribution of image features within each supplied class, \(p(x\mid c)\), together with \(p(c)\), and obtains a classifier by Bayes' rule. Its learned class-conditional model describes features typical of that class.


A discriminative approach learns \(p(c\mid x)\) or a decision function directly. For a linear classifier, the weights describe which feature combinations help separate the supplied classes.


Discriminative training avoids specifying a full feature distribution. A simple generative model can be effective with limited labelled data when its assumptions are appropriate; this is a modelling trade-off, not a guarantee of better accuracy.
